\question[6]
Confronta le tre tecniche degli esempi di laboratorio --- \textbf{override} (Es.\ 1), \textbf{corruption} (Es.\ 2), \textbf{shellcode} (Es.\ 3) --- indicando per ciascuna: cosa si sovrascrive, cosa si ottiene, le precondizioni e la difesa principale.
\begin{solution}
\begin{center}\small
\begin{tabular}{@{}llll@{}}
\toprule
 & \textbf{override} & \textbf{corruption} & \textbf{shellcode}\\
\midrule
Sovrascrive & variabile adiacente & return address & return address + stack\\
Ottiene & dato controllato & salto a funzione esistente & esecuzione codice nuovo\\
Precondizioni & buffer+var adiacenti & offset noto & offset + stack eseguibile\\
Difesa & controllo dei limiti & canary / ASLR & NX / ASLR / canary\\
\bottomrule
\end{tabular}
\end{center}
Dalla corruzione di un dato all'esecuzione di codice arbitrario: complessità e impatto crescenti.
\end{solution}
